\section{Requirements and Success Criteria}

The midterm assessment requires a clear definition of the project requirements and measurable success criteria. The team divides the system into three implementation areas: data structure, frontend, and image algorithm. The requirements below should be treated as the working baseline for the midterm submission.

\subsection{Functional Requirements}

\begin{longtable}{|p{0.14\linewidth}|p{0.45\linewidth}|p{0.31\linewidth}|}
\hline
\textbf{Module} & \textbf{Functional requirement} & \textbf{Success criterion} \\
\hline
Data structure & Define database tables, simulated records, and module interface fields for seat status, users, logs, and administrative actions. & SQL scripts can build the required tables; simulation scripts generate records; fields align with other modules. \\
\hline
Frontend & Present seat status, basic user workflows, and visual feedback for available, occupied, and suspected seats. & Implemented screens or prototypes can show core workflows with sample data or backend API responses. \\
\hline
Image algorithm & Convert camera images or sampled video frames into seat-level states using object detection, ROI matching, and a state machine. & The module outputs \texttt{seat\_id}, \texttt{has\_person}, \texttt{has\_object}, \texttt{status}, and \texttt{suspect\_duration}; it can write CSV and is ready to write Redis/MySQL when local services are configured. \\
\hline
Integration & Use shared fields and repository documentation so modules can be connected later. & README files describe dependencies, commands, outputs, and known limitations. \\
\hline
\end{longtable}

\subsection{Non-Functional Requirements}

\begin{itemize}
  \item \textbf{Reproducibility:} teammates should be able to clone the repository and follow documented setup steps.
  \item \textbf{Modularity:} each team member's work should remain understandable and testable as a separate module.
  \item \textbf{Security and privacy:} no passwords, API keys, personal data, or raw confidential library data should be committed.
  \item \textbf{CPU compatibility for midterm:} the image algorithm baseline should run on a normal laptop CPU.
  \item \textbf{Evidence-based progress:} claims should be supported by code, screenshots, tests, sample outputs, or recorded demonstrations.
\end{itemize}

\subsection{Page Budget for the Full Report}

The requirement document recommends a concise report of approximately 8--12 pages, excluding appendices. A practical page allocation for this team report is:

\begin{itemize}
  \item project overview, requirements, and architecture: about 2--3 pages;
  \item data structure module: about 2 pages after teammates fill the TODO sections;
  \item frontend module: about 2 pages after teammates fill the TODO sections;
  \item image algorithm module: about 2--3 pages, already drafted in this template;
  \item risks, final plan, contribution, and AI-use disclosure: about 1--2 pages plus appendix tables.
\end{itemize}
